\IfFileExists{stacks-project.cls}{%
\documentclass{stacks-project}
}{%
\documentclass{amsart}
}

% For dealing with references we use the comment environment
\usepackage{verbatim}
\newenvironment{reference}{\comment}{\endcomment}
%\newenvironment{reference}{}{}
\newenvironment{slogan}{\comment}{\endcomment}
\newenvironment{history}{\comment}{\endcomment}

% For commutative diagrams we use Xy-pic
\usepackage[all]{xy}

% We use 2cell for 2-commutative diagrams.
\xyoption{2cell}
\UseAllTwocells

% We use multicol for the list of chapters between chapters
\usepackage{multicol}

% This is generally recommended for better output
\usepackage{lmodern}
\usepackage[T1]{fontenc}

% For cross-file-references
\usepackage{xr-hyper}

% Package for hypertext links:
\usepackage{hyperref}

% For any local file, say "hello.tex" you want to link to please
% use \externaldocument[hello-]{hello}
\externaldocument[introduction-]{introduction}
\externaldocument[conventions-]{conventions}
\externaldocument[sets-]{sets}
\externaldocument[categories-]{categories}
\externaldocument[topology-]{topology}
\externaldocument[sheaves-]{sheaves}
\externaldocument[sites-]{sites}
\externaldocument[stacks-]{stacks}
\externaldocument[fields-]{fields}
\externaldocument[algebra-]{algebra}
\externaldocument[brauer-]{brauer}
\externaldocument[homology-]{homology}
\externaldocument[derived-]{derived}
\externaldocument[simplicial-]{simplicial}
\externaldocument[more-algebra-]{more-algebra}
\externaldocument[smoothing-]{smoothing}
\externaldocument[modules-]{modules}
\externaldocument[sites-modules-]{sites-modules}
\externaldocument[injectives-]{injectives}
\externaldocument[cohomology-]{cohomology}
\externaldocument[sites-cohomology-]{sites-cohomology}
\externaldocument[dga-]{dga}
\externaldocument[dpa-]{dpa}
\externaldocument[sdga-]{sdga}
\externaldocument[hypercovering-]{hypercovering}
\externaldocument[schemes-]{schemes}
\externaldocument[constructions-]{constructions}
\externaldocument[properties-]{properties}
\externaldocument[morphisms-]{morphisms}
\externaldocument[coherent-]{coherent}
\externaldocument[divisors-]{divisors}
\externaldocument[limits-]{limits}
\externaldocument[varieties-]{varieties}
\externaldocument[topologies-]{topologies}
\externaldocument[descent-]{descent}
\externaldocument[perfect-]{perfect}
\externaldocument[more-morphisms-]{more-morphisms}
\externaldocument[flat-]{flat}
\externaldocument[groupoids-]{groupoids}
\externaldocument[more-groupoids-]{more-groupoids}
\externaldocument[etale-]{etale}
\externaldocument[chow-]{chow}
\externaldocument[intersection-]{intersection}
\externaldocument[pic-]{pic}
\externaldocument[weil-]{weil}
\externaldocument[adequate-]{adequate}
\externaldocument[dualizing-]{dualizing}
\externaldocument[duality-]{duality}
\externaldocument[discriminant-]{discriminant}
\externaldocument[derham-]{derham}
\externaldocument[local-cohomology-]{local-cohomology}
\externaldocument[algebraization-]{algebraization}
\externaldocument[curves-]{curves}
\externaldocument[resolve-]{resolve}
\externaldocument[models-]{models}
\externaldocument[functors-]{functors}
\externaldocument[equiv-]{equiv}
\externaldocument[pione-]{pione}
\externaldocument[etale-cohomology-]{etale-cohomology}
\externaldocument[proetale-]{proetale}
\externaldocument[relative-cycles-]{relative-cycles}
\externaldocument[more-etale-]{more-etale}
\externaldocument[trace-]{trace}
\externaldocument[crystalline-]{crystalline}
\externaldocument[spaces-]{spaces}
\externaldocument[spaces-properties-]{spaces-properties}
\externaldocument[spaces-morphisms-]{spaces-morphisms}
\externaldocument[decent-spaces-]{decent-spaces}
\externaldocument[spaces-cohomology-]{spaces-cohomology}
\externaldocument[spaces-limits-]{spaces-limits}
\externaldocument[spaces-divisors-]{spaces-divisors}
\externaldocument[spaces-over-fields-]{spaces-over-fields}
\externaldocument[spaces-topologies-]{spaces-topologies}
\externaldocument[spaces-descent-]{spaces-descent}
\externaldocument[spaces-perfect-]{spaces-perfect}
\externaldocument[spaces-more-morphisms-]{spaces-more-morphisms}
\externaldocument[spaces-flat-]{spaces-flat}
\externaldocument[spaces-groupoids-]{spaces-groupoids}
\externaldocument[spaces-more-groupoids-]{spaces-more-groupoids}
\externaldocument[bootstrap-]{bootstrap}
\externaldocument[spaces-pushouts-]{spaces-pushouts}
\externaldocument[spaces-chow-]{spaces-chow}
\externaldocument[groupoids-quotients-]{groupoids-quotients}
\externaldocument[spaces-more-cohomology-]{spaces-more-cohomology}
\externaldocument[spaces-simplicial-]{spaces-simplicial}
\externaldocument[spaces-duality-]{spaces-duality}
\externaldocument[formal-spaces-]{formal-spaces}
\externaldocument[restricted-]{restricted}
\externaldocument[spaces-resolve-]{spaces-resolve}
\externaldocument[formal-defos-]{formal-defos}
\externaldocument[defos-]{defos}
\externaldocument[cotangent-]{cotangent}
\externaldocument[examples-defos-]{examples-defos}
\externaldocument[algebraic-]{algebraic}
\externaldocument[examples-stacks-]{examples-stacks}
\externaldocument[stacks-sheaves-]{stacks-sheaves}
\externaldocument[criteria-]{criteria}
\externaldocument[artin-]{artin}
\externaldocument[quot-]{quot}
\externaldocument[stacks-properties-]{stacks-properties}
\externaldocument[stacks-morphisms-]{stacks-morphisms}
\externaldocument[stacks-limits-]{stacks-limits}
\externaldocument[stacks-cohomology-]{stacks-cohomology}
\externaldocument[stacks-perfect-]{stacks-perfect}
\externaldocument[stacks-introduction-]{stacks-introduction}
\externaldocument[stacks-more-morphisms-]{stacks-more-morphisms}
\externaldocument[stacks-geometry-]{stacks-geometry}
\externaldocument[moduli-]{moduli}
\externaldocument[moduli-curves-]{moduli-curves}
\externaldocument[examples-]{examples}
\externaldocument[exercises-]{exercises}
\externaldocument[guide-]{guide}
\externaldocument[desirables-]{desirables}
\externaldocument[coding-]{coding}
\externaldocument[obsolete-]{obsolete}
\externaldocument[fdl-]{fdl}
\externaldocument[index-]{index}

% Theorem environments.
%
\theoremstyle{plain}
\newtheorem{theorem}[subsection]{Theorem}
\newtheorem{proposition}[subsection]{Proposition}
\newtheorem{lemma}[subsection]{Lemma}

\theoremstyle{definition}
\newtheorem{definition}[subsection]{Definition}
\newtheorem{example}[subsection]{Example}
\newtheorem{exercise}[subsection]{Exercise}
\newtheorem{situation}[subsection]{Situation}

\theoremstyle{remark}
\newtheorem{remark}[subsection]{Remark}
\newtheorem{remarks}[subsection]{Remarks}

\numberwithin{equation}{subsection}

% Macros
%
\def\lim{\mathop{\mathrm{lim}}\nolimits}
\def\colim{\mathop{\mathrm{colim}}\nolimits}
\def\Spec{\mathop{\mathrm{Spec}}}
\def\Hom{\mathop{\mathrm{Hom}}\nolimits}
\def\Ext{\mathop{\mathrm{Ext}}\nolimits}
\def\SheafHom{\mathop{\mathcal{H}\!\mathit{om}}\nolimits}
\def\SheafExt{\mathop{\mathcal{E}\!\mathit{xt}}\nolimits}
\def\Sch{\mathit{Sch}}
\def\Mor{\mathop{\mathrm{Mor}}\nolimits}
\def\Ob{\mathop{\mathrm{Ob}}\nolimits}
\def\Sh{\mathop{\mathit{Sh}}\nolimits}
\def\NL{\mathop{N\!L}\nolimits}
\def\CH{\mathop{\mathrm{CH}}\nolimits}
\def\proetale{{pro\text{-}\acute{e}tale}}
\def\etale{{\acute{e}tale}}
\def\QCoh{\mathit{QCoh}}
\def\Ker{\mathop{\mathrm{Ker}}}
\def\Im{\mathop{\mathrm{Im}}}
\def\Coker{\mathop{\mathrm{Coker}}}
\def\Coim{\mathop{\mathrm{Coim}}}

% Boxtimes
%
\DeclareMathSymbol{\boxtimes}{\mathbin}{AMSa}{"02}

%
% Macros for moduli stacks/spaces
%
\def\QCohstack{\mathcal{QC}\!\mathit{oh}}
\def\Cohstack{\mathcal{C}\!\mathit{oh}}
\def\Spacesstack{\mathcal{S}\!\mathit{paces}}
\def\Quotfunctor{\mathrm{Quot}}
\def\Hilbfunctor{\mathrm{Hilb}}
\def\Curvesstack{\mathcal{C}\!\mathit{urves}}
\def\Polarizedstack{\mathcal{P}\!\mathit{olarized}}
\def\Complexesstack{\mathcal{C}\!\mathit{omplexes}}
% \Pic is the operator that assigns to X its picard group, usage \Pic(X)
% \Picardstack_{X/B} denotes the Picard stack of X over B
% \Picardfunctor_{X/B} denotes the Picard functor of X over B
\def\Pic{\mathop{\mathrm{Pic}}\nolimits}
\def\Picardstack{\mathcal{P}\!\mathit{ic}}
\def\Picardfunctor{\mathrm{Pic}}
\def\Deformationcategory{\mathcal{D}\!\mathit{ef}}
\usepackage{cleveref}
\crefname{originallemma}{original lemma}{original lemmas}
\Crefname{originallemma}{Original lemma}{Original lemmas}
\makeatletter
\newlabel{algebra-lemma-characterize-formally-unramified}{{00UO}{source}{}{sourcepointer-1}{}}
\newlabel{algebra-lemma-characterize-formally-unramified@cref}{{[originallemma][0][]00UO}{[1][0][]source}{}{}{}}
\makeatother
\crefname{originallemma}{original lemma}{original lemmas}
\Crefname{originallemma}{Original lemma}{Original lemmas}
\makeatletter
\newlabel{algebra-lemma-etale-flat-unramified-finite-presentation}{{08WD}{source}{}{sourcepointer-2}{}}
\newlabel{algebra-lemma-etale-flat-unramified-finite-presentation@cref}{{[originallemma][0][]08WD}{[1][0][]source}{}{}{}}
\makeatother
\crefname{originallemma}{original lemma}{original lemmas}
\Crefname{originallemma}{Original lemma}{Original lemmas}
\makeatletter
\newlabel{algebra-lemma-finite-presentation-independent}{{00R2}{source}{}{sourcepointer-3}{}}
\newlabel{algebra-lemma-finite-presentation-independent@cref}{{[originallemma][0][]00R2}{[1][0][]source}{}{}{}}
\makeatother
\crefname{originallemma}{original lemma}{original lemmas}
\Crefname{originallemma}{Original lemma}{Original lemmas}
\makeatletter
\newlabel{algebra-lemma-formally-unramified-unramified}{{00UU}{source}{}{sourcepointer-4}{}}
\newlabel{algebra-lemma-formally-unramified-unramified@cref}{{[originallemma][0][]00UU}{[1][0][]source}{}{}{}}
\makeatother
\crefname{originaltheorem}{original theorem}{original theorems}
\Crefname{originaltheorem}{Original theorem}{Original theorems}
\makeatletter
\newlabel{theorem-descent}{{08XA}{source}{}{sourcepointer-5}{}}
\newlabel{theorem-descent@cref}{{[originaltheorem][0][]08XA}{[1][0][]source}{}{}{}}
\makeatother
\title{Descent along a faithfully flat ring map of finite type, finite presentation, formal unramifiedness, unramifiedness and etaleness}
\author{The Stacks Project Authors}
\date{stacks-project:descent.tex:local-snapshot-20261012; GNU FDL 1.2 or later}
\begin{document}
\maketitle
\noindent\textit{Editorial note.} All delivered mathematical text is the author's own, copied from the pinned Stacks Project chapter file descent.tex; no statement or proof step was rewritten, reordered or removed.\par\medskip
\noindent\textit{Editorial note.} The proof cites [EGA4] twice, once as an alternative to the Algebra lemma used in part (c) and once as an alternative in part (e). Both citations lie in the delivered proof text and are reproduced verbatim; they are redundant with the same-chapter-external Algebra lemmas that the author names first, so no mathematical content depends on them.\par\medskip
\noindent\textit{Editorial note.} The counted spans carry 4 navigation stubs; no same-chapter context span is required and none was added.\par\medskip
\noindent\textit{Editorial note.} The preamble span points at the chapter's own preamble.tex shipped with the frozen source, used verbatim as staged.\par\medskip
\noindent\textit{Allowed standard prerequisites (editorial).}\begin{itemize}
\item Descent of module properties along a faithfully flat ring map. Same chapter, Theorem theorem-descent, invoked at proof lines 1710-1711. Turns the vanishing of a base changed module into the vanishing of the module itself, which is how parts (c) and (e) conclude.
\item Finite generation of the kernel of a map of polynomial algebras can be checked after faithfully flat base change. Algebra, Lemma algebra-lemma-finite-presentation-independent, invoked at proof lines 1694-1696. Supplies the finite set of elements from which the finitely presented approximation A' of A is built in part (b).
\item The differential criterion for formal unramifiedness and the flat, unramified, finitely presented criterion for etaleness. Algebra, Lemma algebra-lemma-characterize-formally-unramified at proof lines 1706-1708, Algebra, Lemma algebra-lemma-formally-unramified-unramified at proof lines 1714-1717, Algebra, Lemma algebra-lemma-etale-flat-unramified-finite-presentation at proof lines 1720-1722. Convert the properties being descended into conditions on the module of relative differentials, on finite type, and on flatness.
\end{itemize}\medskip
\ifdefined\hypertarget\hypertarget{sourcepointer-1}{}\fi
\noindent\textit{Source reference.} Original lemma 00UO: sources/stacks-chow/tags.txt, source lines 1118--1118. Pointer outside the counted span, invoked once for context; the load-bearing argument of the counted statement is the author's proof reproduced in full in the proof span.\par\medskip
\ifdefined\hypertarget\hypertarget{sourcepointer-2}{}\fi
\noindent\textit{Source reference.} Original lemma 08WD: sources/stacks-chow/tags.txt, source lines 11266--11266. Pointer outside the counted span, invoked once for context; the load-bearing argument of the counted statement is the author's proof reproduced in full in the proof span.\par\medskip
\ifdefined\hypertarget\hypertarget{sourcepointer-3}{}\fi
\noindent\textit{Source reference.} Original lemma 00R2: sources/stacks-chow/tags.txt, source lines 988--988. Pointer outside the counted span, invoked once for context; the load-bearing argument of the counted statement is the author's proof reproduced in full in the proof span.\par\medskip
\ifdefined\hypertarget\hypertarget{sourcepointer-4}{}\fi
\noindent\textit{Source reference.} Original lemma 00UU: sources/stacks-chow/tags.txt, source lines 1124--1124. Pointer outside the counted span, invoked once for context; the load-bearing argument of the counted statement is the author's proof reproduced in full in the proof span.\par\medskip
\ifdefined\hypertarget\hypertarget{sourcepointer-5}{}\fi
\noindent\textit{Source reference.} Original theorem 08XA: sources/stacks-descent/pinned-original-source.tex, source lines 1507--1507. Pointer outside the counted span, invoked once for context; the load-bearing argument of the counted statement is the author's proof reproduced in full in the proof span.\par\medskip
\begin{theorem}
\label{theorem-descend-algebra-properties}
If $A \otimes_R S$ has one of the following properties as an $S$-algebra
\begin{enumerate}
\item[(a)]
of finite type;
\item[(b)]
of finite presentation;
\item[(c)]
formally unramified;
\item[(d)]
unramified;
\item[(e)]
\'etale;
\end{enumerate}
then so does $A$ as an $R$-algebra (and of course conversely).
\end{theorem}

\begin{proof}
To prove (a), choose a finite set $\{x_i\}$ of generators of $A \otimes_R S$
over $S$. Write each $x_i$ as $\sum_j y_{ij} \otimes s_{ij}$ with
$y_{ij} \in A$ and $s_{ij} \in S$. Let $F$ be the polynomial $R$-algebra
on variables $e_{ij}$ and let $F \to M$ be the $R$-algebra map sending
$e_{ij}$ to $y_{ij}$. Then $F \otimes_R S\to A \otimes_R S$ is surjective, so
$\Coker(F \to A) \otimes_R S$ is zero and hence $\Coker(F \to A)$
is zero. This proves (a).

\medskip\noindent
To see (b) assume $A \otimes_R S$ is a finitely presented $S$-algebra.
Then $A$ is finite type over $R$ by (a). Choose a surjection
$R[x_1, \ldots, x_n] \to A$ with kernel $I$.
Then $I \otimes_R S \to S[x_1, \ldots, x_n] \to A \otimes_R S \to 0$ is exact.
By Algebra, Lemma \ref{algebra-lemma-finite-presentation-independent}
the kernel of $S[x_1, \ldots, x_n] \to A \otimes_R S$
is a finitely generated ideal. Thus we can find finitely many elements
$y_1, \ldots, y_t \in I$ such that the images of $y_i \otimes 1$ in
$S[x_1, \ldots, x_n]$ generate the kernel of
$S[x_1, \ldots, x_n] \to A \otimes_R S$.
Let $I' \subset I$ be the ideal generated by $y_1, \ldots, y_t$.
Then $A' = R[x_1, \ldots, x_n]/I'$ is a finitely presented $R$-algebra
with a morphism $A' \to A$ such that $A' \otimes_R S \to A \otimes_R S$
is an isomorphism. Thus $A' \cong A$ as desired.

\medskip\noindent
To prove (c), recall that $A$ is formally unramified over $R$ if and only
if the module of relative differentials $\Omega_{A/R}$ vanishes, see
Algebra, Lemma \ref{algebra-lemma-characterize-formally-unramified} or
\cite[Proposition~17.2.1]{EGA4}.
Since $\Omega_{(A \otimes_R S)/S} = \Omega_{A/R} \otimes_R S$,
the vanishing descends by Theorem \ref{theorem-descent}.

\medskip\noindent
To deduce (d) from the previous cases, recall that $A$ is unramified
over $R$ if and only if $A$ is formally unramified and of finite type
over $R$, see
Algebra, Lemma \ref{algebra-lemma-formally-unramified-unramified}.

\medskip\noindent
To prove (e), recall that by
Algebra, Lemma \ref{algebra-lemma-etale-flat-unramified-finite-presentation}
or \cite[Th\'eor\`eme~17.6.1]{EGA4} the algebra
$A$ is \'etale over $R$ if and only if
$A$ is flat, unramified, and of finite presentation over $R$.
\end{proof}

%% Bibliography retained so that the citation carried by the delivered proof text of
%% S5-des-descend-algebra-properties (theorem-descend-algebra-properties) resolves.
%%
%% The delivered proof cites [EGA4] twice, in parts (c) and (e), as an alternative to the
%% Algebra lemmas the author names first.  The entry below is copied character for character
%% from the pinned Stacks Project bibliography file my.bib (the @BOOK{EGA4} record), which is
%% the file the Stacks Project itself uses to resolve this key.  Nothing was rewritten,
%% reordered or invented, and no citation key was changed.
%%
%% Only this one key is cited by the card, and only this one entry is retained.
\begin{thebibliography}{99}

\bibitem[EGA4]{EGA4}
Grothendieck, Alexander and Dieudonn{\'e}, Jean.
\newblock {\'E}l{\'e}ments de g{\'e}om{\'e}trie alg{\'e}brique {IV}.
\newblock Institute des {H}autes {\'E}tudes {S}cientifiques.,
Publications {M}ath{\'e}matiques, vol. 20, 24, 28, 32, 1964-1967.

\end{thebibliography}
\end{document}
